% Generated draft: requires mathematical and source review.
\documentclass[11pt]{article}
\usepackage[margin=25mm]{geometry}
\usepackage{fontspec}
\setmainfont{Latin Modern Roman}
\usepackage{amsmath,amssymb,mathtools,unicode-math}
\setmathfont{Latin Modern Math}
\usepackage{enumitem,needspace,xurl,hyperref}
\setlength{\parindent}{0pt}
\setlength{\parskip}{5pt}
\setlength{\emergencystretch}{3em}
\newcommand{\N}{\mathbb{N}}
\newcommand{\Z}{\mathbb{Z}}
\newcommand{\Q}{\mathbb{Q}}
\newcommand{\R}{\mathbb{R}}
\newcommand{\C}{\mathbb{C}}
\newcommand{\F}{\mathbb{F}}
\newcommand{\A}{\mathbb{A}}
\newcommand{\PP}{\mathbb{P}}
\newcommand{\E}{\mathbb{E}}
\newcommand{\eps}{\varepsilon}
\newcommand{\dd}{\,\mathrm{d}}
\begin{document}
\section*{UnsolvedMath ID: 30000013}
\textbf{Source problem code:} catalogue entry
\subsection{Definitions and mathematical statement}
Let $X$ be a compact complex manifold of complex dimension $n$, and let $\alpha$ be a smooth real $d$-closed $(1,1)$-form. The reality convention is needed for its Hermitian eigenvalues and positivity. Fix a Hermitian metric $\omega$ and let $X(\alpha,\le1)$ be the locus where $\alpha$ has at most one negative eigenvalue relative to $\omega$. This inertia condition is independent of the chosen metric. Including points with a zero eigenvalue does not alter the integral below, since $\alpha^n$ vanishes there. Set
\[
 I(\alpha)=\int_{X(\alpha,\le1)}\alpha^n.
\]
The integral uses the signed top exterior power, not its absolute value: the index-one contribution has negative sign.

Write $[\alpha]\in H^{1,1}_{\mathrm{BC}}(X,\R)$ for the real Bott--Chern class, so representatives differ by $\sqrt{-1}\,\partial\bar\partial$ of real potentials. A K\"ahler current in this class is a closed real current
\[
 T=\alpha+\sqrt{-1}\,\partial\bar\partial\varphi\ge\varepsilon\omega
 \quad\text{for some }\varepsilon>0,
\]
where $\varphi$ is a suitable quasi-plurisubharmonic potential. The usual transcendental volume is the supremum of the top-degree masses of positive representatives with analytic singularities, integrated on their regular loci; the convention here agrees with $\int_X\alpha^n$ for a K\"ahler form, without an extra factor $1/n!$.

The requested transcendental holomorphic Morse inequality is
\[
 I(\alpha)>0\quad\Longrightarrow\quad
 [\alpha]\text{ contains a K\"ahler current and }
 \operatorname{vol}([\alpha])\ge I(\alpha).
\]
The source states this for compact complex manifolds; no K\"ahler or projective hypothesis on $X$ is added. The literature assessment checked on 21 August 2026 attributes the problem to Demailly's transcendental holomorphic Morse inequalities and reports that weak inequalities, special geometries and partial volume estimates are known while the full general lower bound remains unestablished. This is a supplied dated assessment requiring independent verification.
\subsection{Short English statement}
A positive signed integral over the region where a closed real $(1,1)$-form has at most one negative direction should force its class to be big, with volume at least that integral. The target retains the full compact-complex setting.
\subsection{Sources}
\begin{itemize}
\item \emph{Multiplier Ideal Sheaves in Algebraic and Complex Geometry} (2005), Oberwolfach Reports, DOI 10.4171/owr/2004/21, pp. 1108--1166, \url{https://doi.org/10.4171/owr/2004/21}. Proposed year is supplied as 2004; the citation year is retained as given.
\item Literature assessment checked 21 August 2026: ``weak transcendental Morse inequalities,'' \url{https://arxiv.org/abs/1308.2878}; ``partial solution to Demailly's conjecture,'' \url{https://arxiv.org/abs/2309.04978}; Demailly, ``holomorphic Morse inequalities,'' \url{https://www-fourier.ujf-grenoble.fr/~demailly/manuscripts/asympt_cohomology.pdf}.
\item Local UnsolvedMath registry, permanent ID $30000013$, exact source code \texttt{catalogue entry}, created 25 August 2026 and updated 16 September 2026. The clean statement is labelled corrected and verified by notation normalization; the partially-solved literature status is dated curation.
\end{itemize}
\end{document}
